\chapter{Methodology}
\label{ch:method}

\section{Experimental Design and Scope}
\label{sec:method-scope}

The methodology separates model construction, optimisation, and verification. First, a fixed fire scenario is prepared on a spatial grid. Second, numerical deposition coefficients are calculated for every eligible drop origin. Third, a mixed-integer linear program schedules aircraft movement, drops, and refills. Finally, the extracted schedule is replayed independently and compared with the predicted fire history and reported objectives. Controlled cases supply an additional comparison with solutions derived before optimisation.

Static denotes the absence of fire propagation during the horizon, rather than an absence of time-dependent decisions. Aircraft still move, water loads change, and burning cells can become extinguished. There is no natural burnout, re-ignition, evaporation, runoff, stochastic suppression, or FARSITE execution in this experiment. Initially unburned cells remain unburned, and cells already extinguished at the start remain extinguished. A burning cell changes state only when cumulative deposition reaches its threshold.

\section{Spatial and Temporal Representation}
\label{sec:method-grid}

Let $\mathcal{C}$ be a rectangular grid with $N_r$ rows and $N_c$ columns. Cell width and height are $s_x$ and $s_y$ metres, respectively. The initial state of cell $c$ is $z_{c,0}\in\{0,1,2\}$, representing unburned, burning, and extinguished. The set of initially burning cells is
\begin{equation}
    \mathcal{B}=\{c\in\mathcal{C}:z_{c,0}=1\}.
\end{equation}
Only cells in $\mathcal{B}$ require suppression decision states in the optimisation; the other fire states are constants.

Geographic inputs must use a projected coordinate reference system with metre units and a north-up, unrotated raster transform. Ignition polygons are reprojected into the landscape coordinate system. A cell is initially burning when the fraction of its area intersected by the union of ignition polygons is at least a specified threshold, set to $0.5$ in the standard configuration. Polygon intersections are used rather than a cell-centre inclusion rule. Input checks reject incompatible coordinates, missing shapefile components, invalid polygons, and unusable service locations.

The horizon contains $W$ action intervals of duration $\Delta$ minutes, with state boundaries $t=0,\ldots,W$. An action beginning at boundary $t$ and lasting $R$ intervals occupies $[t,t+R)$ and completes at boundary $t+R$. Every selected action must complete by $W$. Distinguishing intervals from boundaries ensures that actions during the final interval have both a resource cost and a recorded fire-state consequence.

\section{Aircraft States and Movement}
\label{sec:method-aircraft}

Let $\mathcal{A}$ be the aircraft set. Aircraft $a$ has capacity $Q_a$ litres, speed $v_a$ metres per minute, and refill duration $R_a^{\mathrm{f}}$ intervals. Integer variables $r_{a,t}$ and $c_{a,t}$ specify its row and column, while $q_{a,t}$ is its integer water load. All aircraft start at the airfield $p_0$ with full capacity:
\begin{equation}
    (r_{a,0},c_{a,0})=p_0,
    \qquad q_{a,0}=Q_a,
    \qquad 0\leq q_{a,t}\leq Q_a.
\end{equation}
Tankers must return to the airfield at $W$. Scooper return is configurable; it is disabled in the standard landscape configuration and enabled in the controlled validation cases.

The displacement charged during interval $t$ is
\begin{equation}
    \ell_{a,t}
    =s_y|r_{a,t+1}-r_{a,t}|
    +s_x|c_{a,t+1}-c_{a,t}|,
    \qquad \ell_{a,t}\leq v_a\Delta.
    \label{eq:static-movement}
\end{equation}
Gurobi absolute-value constraints represent the two coordinate differences. The model can therefore allow movement across several cells in one interval without enumerating every possible origin--destination cell pair. At sampled boundaries, different aircraft cannot occupy the same cell except at the shared airfield.

Equation~\eqref{eq:static-movement} constrains sampled displacement, not a continuous aircraft trajectory. Paths between boundaries are not reconstructed, and crossing paths between samples are not checked for separation. The model also imposes no turning-radius, altitude, terrain-avoidance, or no-fly-zone constraints. Integer positions have an additional consequence: when the permitted interval distance is smaller than both cell dimensions, an aircraft cannot move to another cell. Unused travel allowance does not accumulate across intervals.

\section{Drop and Refill Scheduling}
\label{sec:method-actions}

Let $d_{a,j,t}$ be binary and equal to one when aircraft $a$ begins a drop at initially burning origin $j\in\mathcal{B}$ at boundary $t$. Each drop discharges $D$ litres and lasts $R^{\mathrm{d}}$ intervals. A selected drop requires the aircraft to remain at $j$ throughout service, including its start and completion boundaries. Drop origins are restricted to cells that are still burning when the action begins; drops are not selected from arbitrary unburned grid locations.

For aircraft $a$, let $\mathcal{S}_a$ contain its eligible refill locations. This is the airfield for a tanker and the configured water-source cells for a scooper. Binary $f_{a,s,t}$ selects a refill beginning at $s\in\mathcal{S}_a$. The aircraft remains at that site throughout its refill duration. Movement cannot take place during either service, and the sum of drop and refill actions active in any one aircraft interval is at most one.

Let $u_{a,s,t}$ be the integer volume supplied by a refill. Its constraints are
\begin{align}
    f_{a,s,t}\leq u_{a,s,t}&\leq Q_a f_{a,s,t},\\
    f_{a,s,t}=1&\ \Longrightarrow\
    u_{a,s,t}=Q_a-q_{a,t}.
\end{align}
The lower bound requires at least one litre for a selected refill, preventing a full aircraft from selecting an otherwise redundant zero-volume service. An unselected refill supplies no water.

For every interval, water conservation is enforced by
\begin{equation}
\begin{split}
    q_{a,t+1}={}&q_{a,t}-D\sum_{j\in\mathcal{B}}d_{a,j,t}\\
    &+\sum_{s\in\mathcal{S}_a}
    \sum_{\tau:\,\tau+R_a^{\mathrm{f}}=t+1}u_{a,s,\tau},
\end{split}
\label{eq:static-water-balance}
\end{equation}
together with
\begin{equation}
    D\sum_{j\in\mathcal{B}}d_{a,j,t}\leq q_{a,t}.
\end{equation}
Only feasible action-start indices are included in these sums. Water is committed when a drop starts, so its deduction appears in the next boundary load; deposition on the fire is credited only at drop completion. Refill water becomes available at its completion boundary. These conventions prevent an aircraft from dropping water that is still being replenished or from restoring capacity merely by waiting.

\section{Numerical Water Deposition}
\label{sec:method-deposition}

The footprint is centred on the selected drop-origin cell. In local along-track and cross-track coordinates $(u,v)$, let $b=kh^{1.5}$ be its maximum width, $l_1$ the length of each tapered end, and $l_2$ its plateau length. The centred longitudinal width is
\begin{equation}
    \lambda(u)=b\min\!\left\{1,
        \max\!\left[0,\frac{l_1+l_2/2-|u|}{l_1}\right]\right\}.
\end{equation}
The hexagon has total length $2l_1+l_2$ and lateral extent $|v|\leq\lambda(u)/2$. Since $\sigma(u)=\lambda(u)/6$, its sides truncate the lateral Gaussian at $\pm3\sigma(u)$.

For points within the hexagon with positive width, the normalised deposition density is
\begin{equation}
    \rho(u,v)=\frac{\lambda(u)}{Z}
    \frac{1}{\sqrt{2\pi}\sigma(u)}
    \exp\!\left[-\frac{v^2}{2\sigma(u)^2}\right],
    \label{eq:static-density}
\end{equation}
with zero density outside it and normalising constant
\begin{equation}
    Z=b(l_1+l_2)\operatorname{erf}\!\left(\frac{3}{\sqrt{2}}\right).
\end{equation}
This normalisation accounts for both the longitudinal profile and lateral truncation. The heading rotates the centred footprint in east--north coordinates and is specified clockwise from north.

For origin $j$, let $H_j$ be its translated and rotated hexagon and $P_c$ the polygon of cell $c$. The credited water coefficient is
\begin{equation}
    F_{j,c}=\left\lfloor
        D\int_{P_c\cap H_j}\rho_j(x,y)\,\mathrm{d}x\,\mathrm{d}y
    \right\rfloor.
    \label{eq:static-deposition}
\end{equation}
The lateral integral is evaluated using the normal cumulative distribution function. Adaptive quadrature evaluates the remaining along-track integral, split at polygon vertices and profile breakpoints. Geometric intersections define the integration domain; numerical integration then determines its mass. Coefficients are computed before model construction and are constants within the optimisation.

Rounding down to whole litres ensures conservative credit, with
\begin{equation}
    F_{j,c}\geq0,
    \qquad\sum_{c\in\mathcal{C}}F_{j,c}\leq D.
\end{equation}
Water outside the raster and sub-litre remainders are not redistributed. Water deposited on an unburned cell is recorded during replay but does not extinguish an initially burning cell elsewhere. Consequently, the volume consumed by the aircraft need not equal the volume useful for suppression.

\section{Deterministic Suppression}
\label{sec:method-suppression}

Let $\alpha$ be a specified suppression demand in litres per square metre. The uniform cell threshold is
\begin{equation}
    \theta=\left\lceil\alpha s_xs_y\right\rceil.
\end{equation}
This is a demonstration parameter, not a suppression law calibrated to fuel, fireline intensity, weather, or aircraft performance. The standard value $\alpha=1$ gives $\theta=900$ litres on a $30\times30$ m cell. Controlled cases use different thresholds to isolate particular accounting behaviours.

For $c\in\mathcal{B}$, cumulative delivered water $L_{c,t}$ evolves as
\begin{equation}
    L_{c,0}=0,
    \qquad
    L_{c,t}=L_{c,t-1}
    +\sum_{a\in\mathcal{A}}\sum_{j\in\mathcal{B}}
    \sum_{\tau:\,\tau+R^{\mathrm{d}}=t}F_{j,c}d_{a,j,\tau}.
\end{equation}
Binary $E_{c,t}$ indicates whether the initially burning cell has been extinguished. The initial condition is $E_{c,0}=0$, and the threshold is imposed in both directions:
\begin{align}
    E_{c,t}=1&\ \Longrightarrow\ L_{c,t}\geq\theta,\\
    E_{c,t}=0&\ \Longrightarrow\ L_{c,t}\leq\theta-1,\\
    E_{c,t}&\geq E_{c,t-1}.
\end{align}
Although cumulative-water variables are continuous in the model, their increments are sums of integer coefficients multiplied by binary decisions. The whole-litre threshold therefore permits an exact logical distinction without an arbitrary epsilon gap. Multiple partial deposits accumulate, extinguishment is absorbing, and insufficient water cannot cause a cell to disappear. A drop may begin at $j$ only if
\begin{equation}
    d_{a,j,t}\leq1-E_{j,t}.
\end{equation}

\section{Objective Hierarchy and Solver Procedure}
\label{sec:method-solver}

The primary objective is
\begin{equation}
    \min B,
    \qquad B=\sum_{t=1}^{W}\sum_{c\in\mathcal{B}}(1-E_{c,t}).
    \label{eq:static-primary}
\end{equation}
The initial boundary is excluded because it is fixed, and the final boundary is included so that terminal suppression matters. Multiplication by $\Delta$ gives burning cell-minutes. This sampled burden is neither the number of newly burned cells nor the time until the final drop. It can improve when only part of the initial fire is extinguished.

After the first solve, let $\widehat B$ be the best incumbent value. The second solve adds $B\leq\widehat B$ and minimises
\begin{equation}
    \min\sum_{a\in\mathcal{A}}\sum_{t=0}^{W-1}\ell_{a,t}.
    \label{eq:static-secondary}
\end{equation}
When the first solve proves optimality, this is a lexicographic minimisation of burning cell-steps and then distance. When it times out, the distance stage preserves the best known primary value but does not establish that value's global optimality. Distance is measured in metres and is not interchangeable with total flight time for a heterogeneous fleet.

The stages share a configured solver-time budget, with $80\%$ initially allocated to the primary stage and the unused remainder available to the secondary stage. The requested relative MIP gap is zero. A parked, fully loaded fleet with no actions supplies a feasible initial schedule. Because suppression is an objective rather than a requirement to extinguish every cell, even an ineffective parked solution can be feasible. Controlled cases use a 30 s budget, seed 42, and one solver thread. The standard landscape configuration allows 500 s and automatic thread selection.

The combined status is \texttt{OPTIMAL} only when both stages finish optimally. Otherwise, a returned incumbent is labelled \texttt{FEASIBLE}. Without an incumbent, objective values are left undefined rather than reported as zero. Solver statuses, the primary lower bound, remaining gaps, model dimensions, and elapsed solve time are stored with the extracted schedule. If infeasibility is established and an output directory is available, an irreducible inconsistent subsystem is also written for diagnosis.

\section{Independent Replay and Reproducibility}
\label{sec:method-replay}

The replay procedure reads the extracted routes and actions without querying the solver's decision variables. It checks complete boundary coverage, integer coordinates and loads, grid limits, starting conditions, return requirements, displacement limits, and sampled occupancy. It also verifies action durations and locations, stationary service, non-overlapping actions, drop volumes, refill quantities, and water conservation at every boundary.

The fire history is then reconstructed by applying deposition events at their completion times. Replayed extinguished cells must match the optimiser's prediction at every boundary, and independently recomputed burning cell-steps and distance must match the reported objectives. This provides a second implementation of schedule accounting. It is not a fully independent physical model: replay and optimisation share the scenario, fleet definitions, and precomputed deposition coefficients. Separate geometry tests and analytically constructed cases are therefore needed to check assumptions that replay alone could reproduce incorrectly.

The standard runner saves scenario metadata, solver results, initial and final georeferenced fire grids, replay arrays, and a visual summary. The controlled validation runner stores the expected answer, scenario parameters, actual result, comparison report, model export, solver log, and replay arrays for each case. Unique run directories preserve previous experiments. Drop variables scale with aircraft count, initially burning origins, and horizon length, while pairwise sampled occupancy grows with the square of aircraft count. Avoiding all-grid-cell route-pair enumeration reduces construction overhead, but does not remove the combinatorial difficulty of the scheduling problem.
